% fig81-pick.tex -- Figure 15.x: Pick's theorem on polygon (1,1)-(5,3)-(3,5)-(1,4): A=9, B=8, I=6 (A = I + B/2 - 1).
% fstyle.tex -- shared TikZ styles for the figures of this book (palette from lifeos.sty).
\tikzset{%
  cell/.style={draw=ink, minimum width=7mm, minimum height=6.5mm, inner sep=0pt, font=\footnotesize, fill=white},
  cellh/.style={cell, fill=accentsoft, draw=accent},
  cella/.style={cell, fill=amberbg, draw=amber},
  cellg/.style={cell, fill=goodbg, draw=good},
  cellx/.style={cell, fill=graybg, draw=slate, text=slate},
  vtx/.style={circle, draw=ink, minimum size=6.5mm, inner sep=0pt, font=\footnotesize, fill=white},
  vtxh/.style={vtx, draw=accent, fill=accentsoft},
  vtxa/.style={vtx, draw=amber, fill=amberbg},
  vtxg/.style={vtx, draw=good, fill=goodbg},
  vtxd/.style={vtx, draw=slate, fill=graybg, text=slate},
  lab/.style={font=\footnotesize, text=slate},
  labA/.style={font=\footnotesize, text=accent},
  labM/.style={font=\footnotesize, text=amber},
  idx/.style={font=\tiny, text=slate},
  arr/.style={-{stealth}, thick, ink},
  arrA/.style={-{stealth}, thick, accent},
  arrM/.style={-{stealth}, thick, amber},
  arrG/.style={-{stealth}, thick, good},
  edge/.style={thick, ink},
  edgeA/.style={thick, accent},
  edgeM/.style={thick, amber},
}%
% \cellrow{prefix}{x0}{y}{style}{v0,v1,...}: a row of cells named prefix0, prefix1, ... (pitch 0.75cm)
\providecommand{\cellrow}[5]{\foreach \v [count=\i from 0] in {#5} {\node[#4] (#1\i) at (#2+\i*0.75,#3) {\v};}}%
% \cellidx{prefix}{n}{start}: index labels under cells prefix0 .. prefix(n-1)
\providecommand{\cellidx}[3]{\foreach \i in {0,...,\numexpr#2-1\relax} {\node[idx, below=0pt of #1\i] {\the\numexpr\i+#3\relax};}}%
% \fitfig{<tikzpicture>}: shrink the picture to \linewidth only when it would overflow.
\ifdefined\fitbox\else\newsavebox{\fitbox}\fi%
\providecommand{\fitfig}[1]{\sbox{\fitbox}{#1}\ifdim\wd\fitbox>\linewidth\resizebox{\linewidth}{!}{\usebox{\fitbox}}\else\usebox{\fitbox}\fi}%

\begin{tikzpicture}[font=\footnotesize, x=0.9cm, y=0.9cm]
\draw[slate!35] (0,0) grid (6,6);
\fill[accent!10] (1,1)--(5,3)--(3,5)--(1,4)--cycle; \draw[ink, very thick] (1,1)--(5,3)--(3,5)--(1,4)--cycle;
\fill[accent] (2,2) circle (3pt);
\fill[accent] (2,3) circle (3pt);
\fill[accent] (2,4) circle (3pt);
\fill[accent] (3,3) circle (3pt);
\fill[accent] (3,4) circle (3pt);
\fill[accent] (4,3) circle (3pt);
\fill[good] (1,1) circle (2.6pt);
\fill[good] (1,2) circle (2.6pt);
\fill[good] (1,3) circle (2.6pt);
\fill[good] (1,4) circle (2.6pt);
\fill[good] (3,2) circle (2.6pt);
\fill[good] (3,5) circle (2.6pt);
\fill[good] (4,4) circle (2.6pt);
\fill[good] (5,3) circle (2.6pt);
\node[lab, anchor=west, align=left] at (6.8,5.0) {area $A=9$ (shoelace: $2A=18$)};
\node[font=\footnotesize, text=good, anchor=west] at (6.8,4.0) {$\bullet$ boundary points $B=8$};
\node[labA, anchor=west] at (6.8,3.2) {$\bullet$ interior points $I=6$};
\node[labM, anchor=west, align=left] at (6.8,2.0) {$I=A-\dfrac B2+1=9-4+1=6$ \checkmark};
\end{tikzpicture}
